\renewcommand{\Anno}{2025}

\renewcommand{\NumeroProblema}{2}

\renewcommand{\Testo}{
È ben noto che negli USA le date si scrivono mettendo il mese prima del giorno, mentre in Italia si scrivono mettendo il giorno prima del mese. Per esempio, 2/3 in Italia rappresenta il 2 Marzo, mentre negli USA è il 3 Febbraio. Una data $X/Y$ si dice \emph{ambigua} se ha senso sia negli USA sia in Italia.

\begin{enumerate}
\renewcommand{\labelenumi}{(\alph{enumi})}

\item Per partecipare ad un concorso internazionale, i tantissimi concorrenti (tutti nati nel 2006 o 2007) hanno dovuto compilare un form online in cui era richiesta la data di nascita. Il form controllava che la data indicata avesse senso, e dava un messaggio di errore in caso contrario. Il problema è che tutti i concorrenti europei hanno inserito la data usando lo standard italiano, mentre il form faceva il controllo attenendosi allo standard USA.

Determinare la percentuale di concorrenti europei che hanno ricevuto il messaggio di errore.

\item Enrico e Giuseppe sono nati nello stesso anno. Sappiamo che entrambi sono nati in date ambigue, e che Enrico è nato $d$ giorni dopo Giuseppe, con $74\leq d\leq 78$.

Determinare in che giorno dell'anno è nato Giuseppe.

\item Una data $X/Y$ ambigua si dice \emph{bilanciata} se $X<Y$ e la data corrispondente cade nello stesso giorno della settimana sia negli USA sia in Italia. Ad esempio, quest'anno 2/8 è una data bilanciata, perché sia il 2 Agosto, sia l'8 Febbraio sono caduti di sabato.

Indicato con $B(a)$ il numero di date bilanciate in un certo anno $a$, determinare \emph{quanti} valori distinti ha assunto $B(a)$, al variare di $a$ dal 2001 fino a quest'anno compreso.

[Si ricorda che tutti i 12 mesi dell'anno hanno 31 giorni, tranne Aprile, Giugno, Settembre, Novembre che hanno 30 giorni, e Febbraio che ha 28 giorni, oppure 29 negli anni bisestili.]

\end{enumerate}
}

\renewcommand{\Risposta}{}

\renewcommand{\Aiutino}{}

\renewcommand{\SoluzioneUfficiale}{}

\renewcommand{\SoluzioniAlternative}{}

\renewcommand{\Osservazioni}{}

\renewcommand{\Commenti}{}
